\scanpage{4}

l'ens d'ens ordonné avec \deleted{l'ens} une \deleted{\unsure{sous-ens|linéaire}} \inserted{partie} de $\mathfrak P(F)$ ordonnée par inclusion -- partie stable par $\operatorname{Sup}$ de deux él. majorés ... Par tt somt passent au moins trois arêtes et trois faces.

Notons aussi :

\uline{Cor} Toute face est $\operatorname{Sup}$ des \inserted{ses} arêtes -- tt sommet est $\operatorname{Inf}$ de ses arêtes.

Comme \deleted{tt} sommet d'une face est \unsure{majoré|incidé} par une arête de cette face \inserted{(par constr. de $R''$)}, on voit que

\[
 \operatorname*{Sup}_{a<f}a=\operatorname*{Sup}_{a<f}\operatorname*{Sup}_{s<a}s
 =\operatorname*{Sup}_{s<f}s\overset{\text{cor. 2}}{=}f\qquad\text{cqfd}
\]

(plus dual...)

\uline{Remarque} Dans le cas de l'icosaèdre, sa structure est connue quand on connait $S$ et $A\subset\mathfrak P_2(S)$ -- i.e. on retrouve $F\subset\mathfrak P_3(S)$ à partir de là, comme l'ens des $f\in\mathfrak P_3(S)$ tel que $\deleted{\unsure{\forall}}\mathfrak P_2(f)\unsure{\,|,}\subset A$, \deleted{\unsure{et}} le groupe des autom. de l'icosaèdre s'identifie au groupe des permutations $u$ de $S$ telles que $u(A)=A$ (i.e. $u$ transforme arête en arête -- ou encore, $u$ conserve la ``distance combinatoire'')

\margin{Ce n'est pas vrai pour \uline{tt} polyèdre dont les faces sont des triangles -- p. ex. c'est faux pour l'\uline{icosaèdre gauche}.}
% slanted

% archive sheet number: 3

